% Frittstående fil - kompilerer på egen hånd, er IKKE input i IMAX3012.tex.
%
% "Materiale"-kolonnen lenker til de faktiske delkapitlene i kompendiet
% (IMAX3012.tex/.pdf) via xr-hyper. For at lenkene/numrene skal vise riktig:
%   1) Bygg kompendiet (IMAX3012.tex) FØRST, slik at IMAX3012.aux finnes.
%   2) Bygg deretter denne filen (kjør pdflatex/latexmk to ganger).
%   3) Hold IMAX3012.pdf og fremdriftsplan.pdf i samme mappe -- da fungerer
%      lenkene som hopper inn i kompendium-PDF-en når man klikker på dem.
% Endres kapittelstrukturen i kompendiet må begge dokumentene bygges på nytt.

\documentclass[10pt,a4paper]{article}
\usepackage[a4paper,margin=1cm]{geometry}
\usepackage[norsk]{babel}
\usepackage[utf8]{inputenc}
\usepackage[T1]{fontenc}
\usepackage{lmodern}
\usepackage{amssymb,amsmath,amsfonts}
\usepackage{array}
\usepackage{xcolor}
\usepackage{colortbl}
\usepackage{longtable}
\usepackage{multirow}
\usepackage{enumitem}
\usepackage{caption}
\usepackage{xr-hyper}
\usepackage{hyperref}

\definecolor{fpHeader}{HTML}{8FC1A8}
\definecolor{fpModul1}{HTML}{FDF6DC}
\definecolor{fpModul2}{HTML}{F1E4E1}
\definecolor{fpModul3}{HTML}{D2EEDC}
\definecolor{fpRep}{HTML}{F2F2F2}
\definecolor{fpMatBlue}{HTML}{3B6FA0}
\definecolor{fpBorder}{HTML}{9FAFA8}

\hypersetup{
  colorlinks = true,
  linkcolor = fpMatBlue,
  urlcolor = fpMatBlue,
  citecolor = fpMatBlue,
}

% xr-hyper laster inn IMAX3012.aux råtekst, som inneholder lagrede
% figurtekster/seksjonstitler som bruker kompendiets egendefinerte
% matte-kommandoer (fra style.sty). Disse må derfor også finnes her, ellers
% feiler innlastingen med "Undefined control sequence". Vi trenger ikke at
% de ser nøyaktig riktige ut siden denne teksten aldri vises i dette
% dokumentet - de må bare kunne kompilere.
\providecommand{\cA}{\ensuremath{\mathcal{A}}}
\providecommand{\cB}{\ensuremath{\mathcal{B}}}
\providecommand{\cC}{\ensuremath{\mathcal{C}}}
\providecommand{\cD}{\ensuremath{\mathcal{D}}}
\providecommand{\cE}{\ensuremath{\mathcal{E}}}
\providecommand{\cF}{\ensuremath{\mathcal{F}}}
\providecommand{\cG}{\ensuremath{\mathcal{G}}}
\providecommand{\cH}{\ensuremath{\mathcal{H}}}
\providecommand{\cI}{\ensuremath{\mathcal{I}}}
\providecommand{\cJ}{\ensuremath{\mathcal{J}}}
\providecommand{\cK}{\ensuremath{\mathcal{K}}}
\providecommand{\cL}{\ensuremath{\mathcal{L}}}
\providecommand{\cM}{\ensuremath{\mathcal{M}}}
\providecommand{\cN}{\ensuremath{\mathcal{N}}}
\providecommand{\cO}{\ensuremath{\mathcal{O}}}
\providecommand{\cP}{\ensuremath{\mathcal{P}}}
\providecommand{\cQ}{\ensuremath{\mathcal{Q}}}
\providecommand{\cR}{\ensuremath{\mathcal{R}}}
\providecommand{\cS}{\ensuremath{\mathcal{S}}}
\providecommand{\cT}{\ensuremath{\mathcal{T}}}
\providecommand{\cU}{\ensuremath{\mathcal{U}}}
\providecommand{\cV}{\ensuremath{\mathcal{V}}}
\providecommand{\cW}{\ensuremath{\mathcal{W}}}
\providecommand{\cX}{\ensuremath{\mathcal{X}}}
\providecommand{\cY}{\ensuremath{\mathcal{Y}}}
\providecommand{\cZ}{\ensuremath{\mathcal{Z}}}
\providecommand{\Lf}{\ensuremath{\mathbf{L}}}
\providecommand{\bX}{\ensuremath{\mathbf{X}}}
\providecommand{\bY}{\ensuremath{\mathbf{Y}}}
\providecommand{\LB}[1][\cdot \hspace{1pt} , \cdot]{[\hspace{1pt} #1 \hspace{1pt} ]}
\providecommand{\SSS}{\ensuremath{\mathbb{S}}}
\providecommand{\dd}{\mathop{}\!\mathrm{d}}
\providecommand{\fd}{\mathop{}\!\mathbf{d}}
\providecommand{\E}{\ensuremath{\mathbb{E}}}
\providecommand{\Q}{\ensuremath{\mathbb{Q}}}
\providecommand{\bC}{\ensuremath{\mathbb{C}}}
\providecommand{\B}{\ensuremath{\mathbb{B}}}
\providecommand{\R}{\ensuremath{\mathbb{R}}}
\providecommand{\F}{\ensuremath{\mathbb{F}}}
\providecommand{\N}{\ensuremath{\mathbb{N}}}
\providecommand{\Z}{\ensuremath{\mathbb{Z}}}
\providecommand{\D}{\ensuremath{\mathbb{D}}}
\providecommand{\K}{\ensuremath{\mathbb{K}}}
\providecommand{\vect}[1]{\ensuremath{\mathbf{#1}}}
\DeclareMathOperator{\id}{id}

% Ekstern-referanse til kompendiet (samme mappe). Alle \label i IMAX3012.tex
% blir tilgjengelige her med prefiks "cp-", f.eks. \ref{cp-sect:polarkoord}.
\externaldocument[cp-]{IMAX3012}

% Lenket, dynamisk delkapittel-nummer fra kompendiet.
\newcommand{\fpSec}[1]{\hyperref[cp-#1]{\ref{cp-#1}}}
\newcommand{\fpMat}[1]{#1}

\pagestyle{empty}

\begin{document}

\begin{center}
{\Large\bfseries Fremdriftsplan}\\[.3em]
{\large IMAx3012 -- Matematikk 3B}
\end{center}
\vspace{1em}

Tabellen under viser fremdriftsplanen for kurset, uke for uke, med tilhørende tema, arbeidskrav og relevant materiale fra kompendiet. De tre modulene er fargekodet. Referansene i "Materiale"-kolonnen er klikkbare lenker til riktig delkapittel i kompendiet.

\begingroup
\footnotesize
\setlength{\tabcolsep}{3pt}
\renewcommand{\arraystretch}{1.15}
\arrayrulecolor{fpBorder}

\begin{longtable}{
  |>{\centering\arraybackslash}m{3.0cm}
  |>{\centering\arraybackslash}m{1.0cm}
  |>{\raggedright\arraybackslash}m{6.5cm}
  |>{\centering\arraybackslash}m{2.7cm}
  |>{\raggedright\arraybackslash}m{4.4cm}|
}
\hline
\rowcolor{fpHeader}
\textbf{Modul} & \textbf{Uke} & \textbf{Tema} & \textbf{Relevant arbeidskrav} & \textbf{Materiale} \\
\hline
\endfirsthead
\hline
\rowcolor{fpHeader}
\textbf{Modul} & \textbf{Uke} & \textbf{Tema} & \textbf{Relevant arbeidskrav} & \textbf{Materiale} \\
\hline
\endhead

% ===================== MODUL 1: Multippel-integrasjon =====================
\cellcolor{fpModul1}\multirow{5}{3.0cm}{\centering\textbf{Modul 1}\\[2pt]\textit{Multippel-integrasjon}}
 & \cellcolor{fpModul1}34 &
 \cellcolor{fpModul1}
 \begin{itemize}[nosep,leftmargin=1em]
   \item \textbf{Forelesning 1}
   \begin{itemize}[nosep,leftmargin=1em]
     \item Introduksjon og repetisjon.
     \item Python som verktøy
   \end{itemize}
   \item \textbf{Forelesning 2}
   \begin{itemize}[nosep,leftmargin=1em]
     \item Dobbelintegral i kartesiske koordinater
     \item Iterasjon som metode
     \item Anvendelser
   \end{itemize}
 \end{itemize}
 & \cellcolor{fpModul1}\multirow{2}{2.7cm}{\centering Øving 1} & \cellcolor{fpModul1}\fpMat{Kompendium \hyperref[cp-chap:Intro]{kapittel 0} og kap.\ \fpSec{sect:dobbelint}} \\
\cline{2-3}\cline{5-5}

\cellcolor{fpModul1} & \cellcolor{fpModul1}35 &
 \cellcolor{fpModul1}
 \begin{itemize}[nosep,leftmargin=1em]
   \item \textbf{Forelesning 3}
   \begin{itemize}[nosep,leftmargin=1em]
     \item Dobbeltintegraler i polarkoordinater
   \end{itemize}
   \item \textbf{Forelesning 4}
   \begin{itemize}[nosep,leftmargin=1em]
     \item Python
   \end{itemize}
 \end{itemize}
 & \cellcolor{fpModul1} & \cellcolor{fpModul1}\fpMat{Komp.\ kap.\ \fpSec{sect:polarkoord}-\fpSec{sect:numdobbelt}} \\
\cline{2-5}

\cellcolor{fpModul1} & \cellcolor{fpModul1}36 &
 \cellcolor{fpModul1}
 \begin{itemize}[nosep,leftmargin=1em]
   \item \textbf{Forelesning 5}
   \begin{itemize}[nosep,leftmargin=1em]
     \item Trippelintegral i kartesiske koordinater
     \item Iterasjon som metode
   \end{itemize}
   \item \textbf{Forelesning 6}
   \begin{itemize}[nosep,leftmargin=1em]
     \item Trippelintegraler i sylinderkoordinater
     \item Trippelintegraler i kulekoordinater
   \end{itemize}
 \end{itemize}
 & \cellcolor{fpModul1}\multirow{2}{2.7cm}{\centering Øving 2} & \cellcolor{fpModul1}\fpMat{Komp.\ kap.\ \fpSec{sect:trippelint}-\fpSec{sect:trippelint_andre}} \\
\cline{2-3}\cline{5-5}

\cellcolor{fpModul1} & \cellcolor{fpModul1}37 &
 \cellcolor{fpModul1}
 \begin{itemize}[nosep,leftmargin=1em]
   \item \textbf{Forelesning 7}
   \begin{itemize}[nosep,leftmargin=1em]
     \item Variabelbytte generelt (Jacobi)
   \end{itemize}
   \item \textbf{Forelesning 8}
   \begin{itemize}[nosep,leftmargin=1em]
     \item Python, / anvendelser / oppsamling
   \end{itemize}
 \end{itemize}
 & \cellcolor{fpModul1} & \cellcolor{fpModul1}\fpMat{Komp.\ kap.\ \fpSec{sect:trippelint_andre}-\fpSec{sect:num_met1}} \\
\cline{2-5}

\cellcolor{fpModul1} & \cellcolor{fpModul1}38 &
 \cellcolor{fpModul1}
 \begin{itemize}[nosep,leftmargin=1em]
   \item \textbf{Forelesning 9}
   \begin{itemize}[nosep,leftmargin=1em]
     \item Parametrisering av kurver
     \item vektor-valuerte funksjoner
     \item buelengde-integralet
     \item Linjeintegral for funksjoner
   \end{itemize}
   \item \textbf{Forelesning 10}
   \begin{itemize}[nosep,leftmargin=1em]
     \item parameterflater
     \item Flateintegral
   \end{itemize}
 \end{itemize}
 & \cellcolor{fpModul1} & \cellcolor{fpModul1}\fpMat{Komp.\ kap.\ \fpSec{sect:linint1}-\fpSec{sect:flateint}} \\
\hline

% ===================== MODUL 2: Vektoranalyse =====================
\cellcolor{fpModul2}\multirow{4}{3.0cm}{\centering\textbf{Modul 2}\\[2pt]\textit{Vektoranalyse}}
 & \cellcolor{fpModul2}39 &
 \cellcolor{fpModul2}
 \begin{itemize}[nosep,leftmargin=1em]
   \item \textbf{Forelesning 11}
   \begin{itemize}[nosep,leftmargin=1em]
     \item Vektorfelt
     \item Gradient, divergens og rotasjon
   \end{itemize}
   \item \textbf{Forelesning 12}
   \begin{itemize}[nosep,leftmargin=1em]
     \item Konservative vektorfelt
     \item Potensialfunksjoner
   \end{itemize}
 \end{itemize}
 & \cellcolor{fpModul2}Øving 3 & \cellcolor{fpModul2}\fpMat{Komp.\ kap\ \fpSec{sect:vektorfelt}, \fpSec{div_curl} og \fpSec{sect:konservativ}} \\
\cline{2-5}

\cellcolor{fpModul2} & \cellcolor{fpModul2}40 &
 \cellcolor{fpModul2}
 \begin{itemize}[nosep,leftmargin=1em]
   \item \textbf{Forelesning 13}
   \begin{itemize}[nosep,leftmargin=1em]
     \item Linjeintegraler i vektorfelt
     \item Arbeid og sirkulasjon
   \end{itemize}
   \item \textbf{Forelesning 14}
   \begin{itemize}[nosep,leftmargin=1em]
     \item Fluksintegral i 2D og 3D
   \end{itemize}
 \end{itemize}
 & \cellcolor{fpModul2}\multirow{2}{2.7cm}{\centering Øving 4} & \cellcolor{fpModul2}\fpMat{Komp.\ kap\ \fpSec{sect:linint2}, \fpSec{sect:fluks2d}, \fpSec{sect:fluks3d}} \\
\cline{2-3}\cline{5-5}

\cellcolor{fpModul2} & \cellcolor{fpModul2}41 &
 \cellcolor{fpModul2}
 \begin{itemize}[nosep,leftmargin=1em]
   \item \textbf{Forelesning 15}
   \begin{itemize}[nosep,leftmargin=1em]
     \item Greens teorem
   \end{itemize}
   \item \textbf{Forelesning 16}
   \begin{itemize}[nosep,leftmargin=1em]
     \item Gauss divergenssetning
   \end{itemize}
 \end{itemize}
 & \cellcolor{fpModul2} & \cellcolor{fpModul2}\fpMat{Komp.\ kap\ \fpSec{sect:green},\fpSec{sect:gauss}} \\
\cline{2-5}

\cellcolor{fpModul2} & \cellcolor{fpModul2}42 &
 \cellcolor{fpModul2}
 \begin{itemize}[nosep,leftmargin=1em]
   \item \textbf{Forelesning 17}
   \begin{itemize}[nosep,leftmargin=1em]
     \item Stokes teorem
   \end{itemize}
   \item \textbf{Forelesning 18}
   \begin{itemize}[nosep,leftmargin=1em]
     \item Repetisjon og oppgaveregning med de store teoremene
   \end{itemize}
 \end{itemize}
 & \cellcolor{fpModul2} & \cellcolor{fpModul2}\fpMat{Komp.\ kap\ \fpSec{sect:stokes}} \\
\hline

% ===================== MODUL 3: PDEer =====================
\cellcolor{fpModul3}\multirow{4}{3.0cm}{\centering\textbf{Modul 3}\\[2pt]\textit{PDEer}}
 & \cellcolor{fpModul3}43 &
 \cellcolor{fpModul3}
 \begin{itemize}[nosep,leftmargin=1em]
   \item \textbf{Forelesning 19}
   \begin{itemize}[nosep,leftmargin=1em]
     \item Introduksjon til PDEer
   \end{itemize}
   \item \textbf{Forelesning 20}
   \begin{itemize}[nosep,leftmargin=1em]
     \item Bevaringslover
   \end{itemize}
 \end{itemize}
 & \cellcolor{fpModul3}Øving 5 & \cellcolor{fpModul3}\fpMat{Komp.\ kap\ \fpSec{sect:introPDE}-\fpSec{sect:bevaringslover}} \\
\cline{2-5}

\cellcolor{fpModul3} & \cellcolor{fpModul3}44 &
 \cellcolor{fpModul3}
 \begin{itemize}[nosep,leftmargin=1em]
   \item \textbf{Forelesning 21}
   \begin{itemize}[nosep,leftmargin=1em]
     \item Numerikk for PDEer: Finite differanser og randverdiproblem i 1D
   \end{itemize}
   \item \textbf{Forelesning 22}
   \begin{itemize}[nosep,leftmargin=1em]
     \item Numerikk for PDEer: 2D
   \end{itemize}
 \end{itemize}
 & \cellcolor{fpModul3}\multirow{2}{2.7cm}{\centering Øving 6} & \cellcolor{fpModul3}\fpMat{Komp.\ kap\ \fpSec{sect:endifferanser}} \\
\cline{2-3}\cline{5-5}

\cellcolor{fpModul3} & \cellcolor{fpModul3}45 &
 \cellcolor{fpModul3}
 \begin{itemize}[nosep,leftmargin=1em]
   \item \textbf{Forelesning 23}
   \begin{itemize}[nosep,leftmargin=1em]
     \item Endelig-volum-metode
   \end{itemize}
   \item \textbf{Forelesning 24}
   \begin{itemize}[nosep,leftmargin=1em]
     \item Endelig-volum-metode
   \end{itemize}
 \end{itemize}
 & \cellcolor{fpModul3} & \cellcolor{fpModul3}\fpMat{Komp.\ kap\ \fpSec{sect:endeligvolum}-\fpSec{sect:evm1d}} \\
\cline{2-5}

\cellcolor{fpModul3} & \cellcolor{fpModul3}46 &
 \cellcolor{fpModul3}
 \begin{itemize}[nosep,leftmargin=1em]
   \item \textbf{Forelesning 25}
   \begin{itemize}[nosep,leftmargin=1em]
     \item Endelig-volum-metode
   \end{itemize}
   \item \textbf{Forelesning 26}
   \begin{itemize}[nosep,leftmargin=1em]
     \item Endelig-volum-metode
   \end{itemize}
 \end{itemize}
 & \cellcolor{fpModul3}\multirow{2}{2.7cm}{\centering Øving 7} & \cellcolor{fpModul3}\fpMat{Komp.\ kap\ \fpSec{sect:evm1d}-\fpSec{EVM2D}} \\
\cline{1-3}\cline{5-5}

% ===================== Repetisjon =====================
\rowcolor{fpRep}
\textit{Repetisjon} & 47 &
 \begin{itemize}[nosep,leftmargin=1em]
   \item Repetisjon
 \end{itemize}
 & & \\
\hline

\end{longtable}
\endgroup

\end{document}
